\documentclass[11pt]{article}
\usepackage[review]{acl}

\usepackage{amsmath,amssymb}
\usepackage{booktabs}
\usepackage{multirow}
\usepackage{graphicx}
\usepackage{xcolor}
\usepackage{enumitem}
\usepackage[expansion=false]{microtype}

% Numbers carried over from the ARR author response. Each must be re-checked
% against a result file before submission; then redefine as identity.
\newcommand{\verify}[1]{{\color{red!70!black}#1}}

\title{Scale Is Not Symmetric: Directional Scale Reliability in Chart Question Answering}

\author{Anonymous ACL Submission}

\begin{document}
\maketitle

\begin{abstract}
Chart question answering models must preserve the numerical scale expressed by axes, labels, and units. We study a sparse but consequential failure in which a model predicts a numerically related answer at the wrong scale, such as \textit{0.61} instead of \textit{61}. Our central finding is that this cannot be repaired by unconditional numeric normalization or a one-direction correction: scale correction is \emph{directional} and \emph{regime-dependent}. We construct ScaleQA-Natural, a mined and author-adjudicated diagnostic set of 520 natural scale-conflict candidates (473 image-evaluable), and add a blinded, construction-independent verification of strict scale errors on the ChartQA test set. Using directional scale-contrastive LoRA adapters, prompting, verifiers, calculators, and selectors as probes on an answer-tuned Qwen2.5-VL-7B, we find that (i) a prefer-large adapter recovers most verified strict test-set scale errors (\verify{14/17}) and gives a small but significant held-out gain, whereas compute-matched continuation controls recover almost none (\verify{1/17}); (ii) the same adapter is unsafe as a global policy, because the correct direction reverses across subtypes; and (iii) models often already generate the correctly scaled answer (blind pass@20 of \verify{79.3\%} versus \verify{52.4\%} single-sample), while label-free selectors recover only a small part of that headroom. Our strongest evidence concerns percentage--decimal, factor-of-100 errors on ChartQA-like data; we scope our claims accordingly and argue for regime-aware scale-reliability evaluation with fix--break accounting.
\end{abstract}

\section{Introduction}

Charts make quantitative relations visible, but they also make scale a moving target. The same visual evidence can be verbalized as a percentage, a decimal, a raw value, or a value in thousands or millions. A chart question answering model must therefore do more than find the relevant bar, line, or label: it must decide which scale the answer should use.

Modern multimodal large language models (MLLMs) achieve strong aggregate accuracy on chart QA benchmarks such as ChartQA \citep{masry2022chartqa}. Aggregate accuracy can nevertheless hide rare failures that are costly when they occur. We study \emph{scale inconsistency}: a model gives a numerically related answer under the wrong scale, such as predicting \textit{0.61} when the chart question requires \textit{61}. This is not a generic numeric miss. The model has often found the right quantity but has failed at the last reliability step: choosing the scale in which to express it. In financial, scientific, or policy charts, an answer that is off by a factor of 100 while looking plausible is more dangerous than an obviously wrong one.

The unconditional interventions we test do not solve this. Normalizing all percentages, training an adapter that prefers larger numbers, or asking the model to self-check units each helps some cases and hurts others. A legitimate answer such as \textit{0.0593} should not be expanded; an axis in millions should not be treated like a percent display; and a prediction that is already too large needs the opposite intervention from one that is too small. The central question is therefore not whether a model can \emph{produce} a scaled alternative, but whether a system can \emph{select} the correct scale regime: the direction, unit context, and output convention required by the chart and question.

Natural strict scale errors are sparse under any single model and split, so the evidence must be read at two levels. Across the ChartQA validation/test outputs we audit, strict factor-of-100 candidates account for 0.12--1.09\% of examples (Table~\ref{tab:strict_prevalence}). We treat the automatic ratio rule as a \emph{candidate filter}, not as proof of a scale error, and verify test-set candidates with a blinded protocol in which the annotator sees only the chart, question, gold answer, and prediction (\S\ref{sec:strict}). To understand the mechanism behind these sparse failures, we also construct ScaleQA-Natural, a mined and author-adjudicated set of 520 natural scale-conflict candidates, 473 of which are directly image-evaluable (ScaleQA-Natural-473). It is a diagnostic set, not a replacement for official ChartQA evaluation.

Our main empirical finding is that scale correction is directional and regime-dependent. In an answer-tuned Qwen2.5-VL-7B setting, a prefer-large scale-contrastive adapter improves held-out ChartQA test accuracy by +0.72 points (paired bootstrap 95\% CI $[+0.38,+1.09]$) and recovers \verify{14/17} verified strict test-set scale errors; compute-matched continuation controls with ordinary answers, random numeric negatives, or non-directional pairs recover at most \verify{1/17}, so the effect is attributable to directional scale supervision rather than extra training. The aggregate gain is small because the failure is rare; the targeted effect is large. Yet the same adapter is unsafe when applied globally: on ScaleQA-Natural it breaks more examples than it fixes, and the best direction reverses across subtypes. A blind-generation analysis shows the model often already produces the correctly scaled answer (pass@20 \verify{79.3\%} versus \verify{52.4\%} single-sample), while label-free selectors recover only a small part of that headroom. Candidate generation and regime selection both matter, and regime selection remains the larger open gap.

\noindent\textbf{Contributions.}
We contribute (i) a verified strict scale-error recovery probe on ChartQA test and ScaleQA-Natural, an audited diagnostic set of natural scale conflicts;
(ii) evidence, with compute-matched controls, that directional scale supervision repairs a verified error slice but is unsafe as a global policy;
(iii) a decomposition of candidate generation from regime selection, using prompts, verifiers, calculators, likelihood scoring, blind sampling, and label-free selectors; and
(iv) a scoped conclusion: for percentage--decimal, factor-of-100 errors on ChartQA-like data, reliable scale-regime selection, not candidate availability, is the central unresolved problem.

\section{Related Work}

\paragraph{Chart question answering.}
Chart question answering requires models to combine visual perception, OCR, structural grounding, and numerical reasoning. Earlier benchmarks such as DVQA, FigureQA-style synthetic chart reasoning, and PlotQA emphasize controlled visual reasoning over chart data \citep{kafle2018dvqa,methani2020plotqa}. ChartQA \citep{masry2022chartqa} evaluates both human-written and machine-generated chart questions that involve value lookup, comparison, and arithmetic. Recent chart-specialized systems add plot-to-table reasoning, chart derendering, chart-specific pretraining, and visual instruction tuning \citep{liu2023deplot,liu2023matcha,masry2023unichart,han2023chartllama,masry2024chartgemma}. Recent benchmarks stress more complex and realistic charts \citep{chartbench2024,wang2024charxiv,masry2025chartqapro,wei2025chartmind}, and ChartInsights \citep{wu2024chartinsights} evaluates low-level chart tasks with fine-grained error taxonomies. These diagnostics slice errors by task type, chart type, or reasoning step; none isolates the case in which the model finds the right quantity but reports it under the wrong scale, nor tracks the \emph{direction} of that error. ScaleQA-Natural is complementary: it is a narrow, directional slice that can be mined from any of these benchmarks, and our fix--break accounting applies to any correction evaluated on them.

\paragraph{Numerical reliability in MLLMs.}
Numerical errors in multimodal models can arise from visual grounding failures, OCR mistakes, wrong axis interpretation, arithmetic errors, or answer normalization artifacts. Scale inconsistency differs from a generic numeric miss because the prediction is close under a simple multiplicative transformation. This makes it easy to miss in aggregate error analysis and risky in deployment, since an answer like 0.61 can look superficially related to 61 while being semantically wrong. The distinction also changes the evaluation target: the question is not only whether a model can compute or retrieve a number, but whether it can choose the output convention supported by the chart's axis, unit, and question wording.

\paragraph{Diagnostic evaluation and correction policies.}
Most chart QA work reports aggregate task accuracy, while reliability failures are often analyzed post hoc. We instead treat scale conflicts as a diagnostic slice with inclusion criteria, an audit trail, and fix--break accounting. A useful correction must recover scale errors, preserve non-scale examples, and avoid applying the wrong direction. The high oracle gaps in our experiments further distinguish candidate generation from candidate selection, a failure mode not captured by single-answer accuracy alone. We use TextVQA \citep{singh2019towards} only as an adjacent text-rich safety check.

\paragraph{Parameter-efficient and preference-style adaptation.}
LoRA updates a model through low-rank trainable matrices while keeping the base weights fixed \citep{hu2021lora}. This makes it attractive for targeted reliability interventions in large MLLMs. Contrastive and preference-style objectives, including direct preference optimization \citep{rafailov2023direct}, optimize relative preferences between candidate outputs. Our objective is more local: for a chart and question, we contrast a preferred numeric answer with a scale-related alternative, and study how the direction of this contrast changes behavior.

\section{Diagnosing Scale Inconsistency}

We use three views of the same reliability problem, because no single metric tells the whole story: official ChartQA accuracy measures task utility, verified strict scale-error recovery measures whether real scale mistakes can be repaired, and ScaleQA-Natural exposes why a repair that helps one regime can damage another.

\paragraph{Strict scale-error candidates.}
\label{sec:strict}
Let $y$ be the gold answer and $\hat{y}$ a model prediction. We extract a single numeric value $n(\cdot)$ when possible and flag an example as a \emph{strict scale-error candidate} if the baseline prediction is incorrect, both answers contain one nonzero number, and their ratio is approximately $100$ or $1/100$:
\begin{equation}
\begin{aligned}
\frac{|n(\hat{y})|}{|n(y)|} &\in [80,125], \\
\text{or}\quad
\frac{|n(\hat{y})|}{|n(y)|} &\in [0.008,0.0125].
\end{aligned}
\end{equation}
A percent sign is stripped before parsing, so \textit{0.61\%} is parsed as $0.61$. This rule is a \emph{candidate filter}: a ratio near 100 can also arise from a wrong entity, a wrong operation, or a coincidental numeric miss. We therefore verify candidates manually before using them as a recovery probe.

\paragraph{Blinded verification on ChartQA test.}
To obtain evidence that is independent of the ScaleQA-Natural construction, we scanned all numeric-answer ChartQA test examples with the ratio rule and verified every strict candidate through a blinded interface that shows only the chart, question, gold answer, and model prediction, hiding the candidate group, prediction/gold ratio, and relation label. Of \verify{18} strict factor-of-100 candidates, \verify{17} were verified as true scale errors and \verify{16} as clean scale errors with no confounding issue; among \verify{30} randomly sampled non-candidates, \verify{1} was judged a missed scale error. The verified set (\verify{17} examples) is the recovery probe used in \S\ref{sec:experiments}. Baseline accuracy on it is zero by construction, so we report recovery as counts with Wilson 95\% intervals, never as a population-level rate. Item identifiers, labels, and the full list of examples are given in Appendix~\ref{app:strict_verified}.

\paragraph{ScaleQA-Natural.}
Because strict errors from a single model and split are sparse, we mine scale-conflict candidates from ChartQA predictions across model outputs and split shards and review them with a machine-assisted annotation workflow. ScaleQA-Natural makes this rare failure mode measurable without turning it into a hidden training set or leaderboard. It contains 520 author-adjudicated positives, of which 473 are directly image-evaluable; we call this subset \emph{ScaleQA-Natural-473} and use it for all image-conditioned experiments. Released items include question, gold answer, source row, model prediction, subtype, direction, review label, and audit metadata; images are referenced through original ChartQA rows. We also retain the 403-item seed set, \emph{ScaleQA-Core}, and a 192-example InternVL-mined subset for source-distribution analysis.


\paragraph{Annotation pipeline and what the audit does and does not establish.}
The pipeline has three stages that we report separately (Table~\ref{tab:annotation_audit}, appendix). (1)~\emph{First-pass labeling}: 739 mined candidates receive include/review/reject labels. (2)~\emph{Author adjudication}: a deterministic metadata audit is compared with the first-pass labels; the $\kappa=0.443$ reported there is agreement between this automatic audit and the first-pass labels, \emph{not} human--human agreement, and it flags 108 of the eventual positives for manual adjudication. (3)~\emph{Inclusion audit}: a second human reviewer audited the 520 final positives and made no changes. This second pass reviewed retained positives only, so it provides no estimate of false negatives in the rejected or uncertain pool, and its 100\% agreement should not be read as independent initial label reliability. We therefore describe ScaleQA-Natural as an author-adjudicated, inclusion-audited diagnostic set, not an independently human-authored benchmark. The blinded ChartQA-test verification above is the construction-independent check on the phenomenon itself.


\paragraph{Split hygiene.}
The main Qwen scale adapters are trained from ChartQA training rows and training-set prediction/log-probability caches. ChartQA validation is used for development; held-out ChartQA test is reserved for the main task result. ScaleQA-Natural is used for diagnosis and policy-selection analysis, not for training the reported Qwen scale adapters. \verify{We checked row-level overlap between ScaleQA-Natural and the ChartQA training rows used for answer tuning and scale-pair training; [state the number of overlapping rows here, and if nonzero, report adapter results on the row-disjoint subset in Appendix~\ref{app:overlap}].} Automatically mined candidates that remain uncertain or are rejected often involve plausible ratios but insufficient evidence that prediction and gold refer to the same chart quantity; these are excluded rather than treated as positives (Table~\ref{tab:negative_audit_examples}).

\paragraph{Annotation guideline.}
Annotators mark an item as a scale diagnostic only when the prediction and gold answer are both grounded in the chart question and differ primarily by a simple multiplicative scale relation, such as percentage versus decimal form or a visible unit-scale mismatch. Items are rejected when the prediction reflects the wrong visual entity, wrong operation, OCR failure, ambiguous wording, duplicate evidence, or an arbitrary numeric miss. Ambiguous cases go to a review queue and are not included without adjudication. Annotation was performed from the chart image together with the question and answers, not from metadata alone.

\paragraph{Why this decomposition matters.}
Overall accuracy can remain stable while scale reliability changes sharply, and a correction can recover a few strict errors while breaking many non-scale examples. ChartQA test asks whether an intervention helps the task, the verified strict subset asks whether it repairs natural scale errors, and ScaleQA-Natural asks why the same intervention helps one regime but harms another. We therefore report full accuracy, strict recovery counts, non-scale accuracy, and fix--break accounting whenever possible.

\section{Diagnostic Probes and Policy Selection}

We freeze the MLLM backbone and train LoRA parameters $\Delta$ \citep{hu2021lora}. For a chart image $x$, question $q$, and answer $y$, the adapted model defines $p_{\theta,\Delta}(y\mid x,q)$ with frozen $\theta$. The methods below form a probe ladder rather than a single proposed fix. Directional adapters test whether scale preference can be moved; compute-matched controls test whether the movement is due to the scale objective; clean-direct controls test whether it is safe without answer tuning; blind sampling tests whether the correct answer is generated at all; and selectors test whether candidate evidence can be turned into a deployable regime decision.

\begin{figure}[t]
\centering
\small
\setlength{\fboxsep}{4pt}
\fbox{\begin{minipage}{0.95\linewidth}
\textbf{Chart + question}
$\rightarrow$
\textbf{Base answer} (e.g., \textit{0.61})\\[-0.15em]
\rule{\linewidth}{0.3pt}\\[-0.2em]
\textbf{Candidate pool:}
base, self-check answer, directional adapter answers, blind samples, explicit scale variants.\\[-0.15em]
\textbf{Regime evidence:}
question form, answer magnitude, percent/unit tokens, axis labels, tick values, candidate likelihoods.\\[-0.15em]
\textbf{Policy decision:}
base, self-check, prefer-large, prefer-small, decimal guard, or unit guard.\\[-0.15em]
\rule{\linewidth}{0.3pt}\\[-0.1em]
\begin{tabular}{@{}l@{}}
\emph{Prediction too small:} $0.61 \rightarrow 61$ \\
\emph{Prediction too large:} $61 \rightarrow 0.61$ \\
\emph{True small decimal:} keep $0.0593$ \\
\emph{Unit mismatch:} millions vs. raw values \\
\end{tabular}
\end{minipage}}
\caption{Regime-aware diagnostic view. The experiments test whether scale-related candidates can be generated and, more importantly, whether a deployable policy can select the scale regime required by the chart and question.}
\label{fig:regime_pipeline}
\end{figure}

\paragraph{Directional adapters as probes.}
For numeric training examples we construct scale-related pairs $(y^+,y^-)$ in which $y^-$ is $y^+$ rescaled by a factor of $100$ or $1/100$ (with percent signs handled as in \S\ref{sec:strict}; full pair-construction rules, rounding, and sign handling are in Appendix~\ref{app:pairs}). \emph{Prefer-large} pairs take the larger-scale value as $y^+$; \emph{prefer-small} pairs take the smaller. Tables retain the training-time labels \emph{large-positive}/\emph{small-positive}, which name the same preferences. The names describe the training preference, not the observed error direction: an answer that is too small requires a prefer-large correction, and vice versa. We also evaluate balanced mixtures, all mined pairs, and real-conflict pairs mined from model disagreements.

\paragraph{Objective.}
We compute mean token negative log-likelihood $\mathrm{NLL}(y\mid x,q)$ and minimize
\begin{equation}
\begin{aligned}
\mathcal{L}_{\mathrm{scale}}
=\mathrm{softplus}\big(&-\beta[
\mathrm{NLL}(y^-) \\
&-\mathrm{NLL}(y^+)-m]\big),
\end{aligned}
\end{equation}
where $m$ is a margin and $\beta$ controls sharpness. The loss encourages the preferred answer to have lower NLL than its scale-related alternative.

\paragraph{Adaptation regimes and controls.}
The main positive result uses an \emph{answer-tuned} starting point: scale adaptation continues from an answer-only LoRA adapter in the Qwen2.5-VL-7B pipeline for 300 additional steps. Because any continuation could help simply through more optimization, we add \emph{compute-matched controls} that continue the same answer-tuned adapter for the same 300 steps with (a)~ordinary answer examples, (b)~random or shuffled numeric negatives, and (c)~non-directional balanced pairs with matched answer frequency. We also test \emph{clean direct} scale-only tuning from the raw backbone as a safety control. At inference, main results use the same deterministic short-answer decoding template as the baseline, without reranking or calculators.

\paragraph{Blind-generation oracle.}
To test whether candidate oracles are high only because the pool contains deterministic scale transforms of the base answer, we also draw $20$ stochastic samples from the answer-tuned model with no scale prompt and no transformation, and report single-sample accuracy against pass@20. This separates ``the model can generate the right scale'' from ``our transformation rules reproduce the mined relation.''

\paragraph{Label-free policy selection.}
\label{sec:router_method}
ScaleQA-Natural suggests that no single direction is safe across all regimes. We therefore evaluate selectors that choose among base answering, self-check prompting, unit/axis guards, and directional adapters, using question text, base prediction text, numeric features, and chart-context strings extracted from the image (title, axis labels, tick values, legends, units). We call a selector \emph{label-free} when it uses no gold subtype or direction field; this is a weaker guarantee than an untouched test set, and we report selector results under a fixed partition and leave-source-out evaluation accordingly (\S\ref{sec:selector}). Oracle subtype, policy-pool, and candidate-pool results are upper bounds whose role is to show how much headroom remains if regime detection were solved.

\section{Experiments}
\label{sec:experiments}

\paragraph{Setup.}
The main experiments use Qwen2.5-VL-7B-Instruct \citep{qwen25vl2025}. We use ChartQA validation (1{,}920 examples) for development and held-out ChartQA test (2{,}500 examples) for the main task result \citep{masry2022chartqa}. We also evaluate InternVL2.5-1B/8B/26B \citep{internvl25} and numeric subsets of PlotQA \citep{methani2020plotqa}, ChartBench \citep{chartbench2024}, MathVista \citep{lu2023mathvista}, and InfographicVQA \citep{mathew2022infographicvqa}. Numeric answers are scored with the repository evaluator: normalized exact match for text answers and 5\% tolerance for numeric answers. The final Qwen scale adapters use $K=300$ pairs, 300 steps, learning rate $10^{-6}$, $\beta=2.0$, and margin $m=0.10$. All accuracy differences are in percentage points.

\subsection{Held-Out ChartQA and Verified Strict Errors}

The first question is deliberately narrow: can a scale-directed update fix natural scale errors without damaging the main task, and is the effect due to the scale objective rather than extra training?

\begin{table}[t]
\centering
\small
\resizebox{\linewidth}{!}{\begin{tabular}{llrrr}
\toprule
System & Starting point & Test acc. & Std. & Delta \\
\midrule
Raw Qwen2.5-VL-7B & Raw backbone & 0.7804 & -- & -- \\
Answer-tuned LoRA & Answer-only & 0.8860 & -- & +0.1056 \\
Answer-tuned + large-positive scale & Answer-tuned & 0.8932 & 0.0003 & +0.0072 \\
Clean direct large-positive scale & Raw backbone & 0.7456 & -- & -0.0348 \\
Clean direct small-positive scale & Raw backbone & 0.7136 & -- & -0.0668 \\
\bottomrule
\end{tabular}
}
\caption{Held-out ChartQA test attribution for Qwen2.5-VL-7B. The scale update is measured on top of the answer-tuned adapter; clean direct scale-only tuning is a safety control. ``Large-positive'' denotes the prefer-large adapter.}
\label{tab:official_test}
\end{table}

Table~\ref{tab:official_test} separates answer tuning from scale tuning. The raw pipeline obtains 78.04 on ChartQA test; answer tuning reaches 88.60. Continuing from that adapter, prefer-large scale-contrastive tuning reaches 89.32 over five seeds, +0.72 points with paired bootstrap 95\% CI $[+0.38,+1.09]$. Clean scale-only adaptation from the raw backbone hurts test accuracy by 3.5--6.7 points. The aggregate gain is small because strict scale errors are rare (Table~\ref{tab:strict_prevalence}); non-scale accuracy moves only from 89.35 to 89.66 (Table~\ref{tab:test_subset_results}), and TextVQA-500 is unchanged at 94.20 (a small adjacent-task check, not a general safety guarantee).

\begin{table}[t]
\centering
\small
\resizebox{\linewidth}{!}{\begin{tabular}{lcc}
\toprule
Continuation from answer-tuned adapter (300 steps, 3 seeds) & ChartQA val.\ acc. & Verified strict test errors recovered \\
\midrule
None (answer-tuned baseline) & \verify{0.8640} & \verify{0/17} \; [0.000, 0.184] \\
Ordinary answer replay & \verify{$\approx$0.8640} & \verify{1/17} \; [0.010, 0.270] \\
Random / shuffled numeric negatives & \verify{$\approx$0.8640} & \verify{[fill]} \\
Non-directional balanced pairs & \verify{$\approx$0.8640} & \verify{[fill]} \\
Directional scale pairs (prefer-large) & \verify{0.8695} & \verify{12/17} \; [0.469, 0.867] \\
Prefer-large adapter, main setting (5 seeds) & -- & \verify{14/17} \; [0.590, 0.938] \\
\bottomrule
\end{tabular}
}
\caption{Compute-matched controls. Every row continues the same answer-tuned adapter for the same number of steps; only the continuation data differs. Recovery is reported as counts over the \verify{17} blind-verified strict test errors with Wilson 95\% intervals.}
\label{tab:controls}
\end{table}

Table~\ref{tab:controls} answers the attribution question. On the verified strict test errors the answer-tuned baseline recovers \verify{0/17}. Continuing with ordinary answer examples recovers \verify{1/17}, and non-directional or random-negative continuation is statistically tied with the baseline on ChartQA validation. Directional scale supervision recovers \verify{12/17} in the three-seed control setting and \verify{14/17} with the main five-seed prefer-large adapter. The intervals are wide, but they exclude the control outcomes: the targeted effect comes from directional scale supervision, not from 300 more steps. This is a verified diagnostic-slice result, not a population-level rate (per-seed counts: Appendix~\ref{app:strict_verified}).

\subsection{ScaleQA-Natural}
\label{sec:scaleqa_results}

The second question is whether the same intervention is safe when scale conflicts are collected across regimes rather than filtered to one verified recovery subset.

\begin{table*}[t]
\centering
\small
\resizebox{\textwidth}{!}{\begin{tabular}{lrrrrrrrrr}
\toprule
Set & N & Raw & Ans. & Prompt & Large & Bal. & All & Real & Small \\
\midrule
ScaleQA-Natural-520 & 473 & 0.4567 & 0.5095 & 0.4419 & 0.4951 & -- & -- & -- & -- \\
ScaleQA-Core & 403 & 0.4591 & 0.4963 & 0.4293 & 0.4844 & 0.5092 & 0.5047 & 0.5002 & 0.4174 \\
InternVL-mined & 192 & 0.1510 & 0.3281 & 0.0990 & 0.1896 & 0.3323 & 0.2771 & 0.2510 & 0.7490 \\
\bottomrule
\end{tabular}
}
\caption{ScaleQA-Natural diagnostic results. The first row is the image-evaluable ScaleQA-Natural-473 subset; ScaleQA-Core and InternVL-mined report five-seed pair-mode means. Large/Small denote prefer-large/prefer-small adapters.}
\label{tab:scaleqa_results}
\end{table*}

On ScaleQA-Natural-473, answer tuning improves the raw backbone from 45.67 to 50.95, a scale-aware prompt drops to 44.19, and the global prefer-large adapter reaches 49.51, \emph{below} the answer-tuned base. On ScaleQA-Core, balanced and all-pair variants outperform the answer-tuned base; on the InternVL-mined subset, prefer-small adaptation is far stronger than every other global policy (74.90 versus 33.23 balanced).

\begin{table*}[t]
\centering
\small
\resizebox{\textwidth}{!}{\begin{tabular}{lrrrrrrr}
\toprule
ScaleQA-Natural-473 subset & $N$ & Answer-tuned & Prefer-large & Self-check & $\Delta$(large$-$ans.) [95\% CI] & Fix & Break \\
\midrule
Items from ChartQA training rows & 363 & 50.96 & 46.89 & 54.27 & $-4.1$ [$-7.8$, $-0.4$] & 4.7 & 8.8 \\
Items from ChartQA val/test rows (row-disjoint) & 110 & 50.91 & 58.18 & 55.45 & $+7.3$ [$+1.5$, $+13.6$] & 9.1 & 2.7 \\
All & 473 & 50.95 & 49.51 & 54.55 & $-1.4$ [$-4.7$, $+1.7$] & 5.7 & 7.4 \\
\midrule
\multicolumn{8}{l}{\emph{By direction} (answer-tuned / prefer-large / self-check):} \\
Prediction too large, training rows & 205 & 30.7 & 15.3 & 45.4 & & & \\
Prediction too small, training rows & 158 & 77.2 & 87.8 & 65.8 & & & \\
Prediction too large, row-disjoint & 36 & 22.2 & 14.4 & 33.3 & & & \\
Prediction too small, row-disjoint & 74 & 64.9 & 79.5 & 66.2 & & & \\
\bottomrule
\end{tabular}
}
\caption{Row-overlap analysis on ScaleQA-Natural-473. Prefer-large is the five-seed mean; $\Delta$ is the paired per-item difference with a bootstrap 95\% CI; Fix/Break are the percentages of items the adapter corrects/breaks relative to the answer-tuned model. The direction asymmetry is the same on both subsets; the aggregate sign differs because 56\% of training-row items are prediction-too-large versus 33\% of row-disjoint items.}
\label{tab:row_overlap}
\end{table*}

Table~\ref{tab:row_overlap} separates the items by whether their ChartQA row was seen during answer tuning. The aggregate ``adapter hurts'' result is carried entirely by the training-row items ($-4.1$ points, CI $[-7.8,-0.4]$, breaking 8.8\%). On the 110 row-disjoint items the same adapter \emph{helps} ($+7.3$ points, CI $[+1.5,+13.6]$, fixing 9.1\% and breaking 2.7\%) and edges self-check (58.18 versus 55.45). The direction breakdown shows why: in both subsets prefer-large helps when the prediction is too small and hurts when it is already too large; the subsets differ mainly in their mix of directions. Directionality is the stable finding; the sign of a global adapter's aggregate effect on a mined set depends on the set's direction composition, itself a reason not to deploy unconditional correction.

Training-free prompting baselines (Table~\ref{tab:tta_baselines_app}, appendix) are stronger than the global adapter on the full set. A self-check prompt reaches 54.55; majority voting adds little, while oracle prompt choice reaches 67.44. The question becomes whether the right answer is missing or merely not selected.

\begin{table}[t]
\centering
\small
\resizebox{\linewidth}{!}{\begin{tabular}{lc}
\toprule
Candidate source on ScaleQA-Natural-473 & Accuracy / coverage \\
\midrule
Single greedy sample (answer-tuned) & \verify{0.524} \\
Blind pass@20, no scale prompt, no transforms & \verify{0.793} \\
Prompt + adapter candidates (policy-pool oracle) & 0.9004 \\
Top-8 candidates incl.\ deterministic scale variants & 0.9175 \\
Augmented structured oracle & 0.9260 \\
\bottomrule
\end{tabular}
}
\caption{Candidate coverage on ScaleQA-Natural-473 decomposed by candidate source. Blind sampling uses no scale prompt and no deterministic transformation, so its coverage is not implied by the mining rule.}
\label{tab:blind_oracle}
\end{table}

Table~\ref{tab:blind_oracle} addresses a construction concern directly: because ScaleQA-Natural is mined through multiplicative relations and the structured oracle pools contain explicit scale variants, oracles over those pools are high partly by construction. The blind-generation row removes both. With no scale prompt and no transformation, \verify{20} stochastic samples already contain a correct answer for \verify{79.3\%} of examples, against \verify{52.4\%} for a single sample. Deterministic variants raise coverage only further, to 91.75--92.60; most of the gap to deployable methods is present without them.

\subsection{Cross-Backbone and External Checks}

The final question is scope: does the phenomenon extend beyond the main Qwen setting, and does the intervention transfer with it? Table~\ref{tab:transfer_boundary} (appendix) summarizes the picture. InternVL2.5-1B and 8B can be adapted with the same scale-pair recipe, and InternVL2.5-26B serves as a stronger diagnostic backbone on ScaleQA-Natural, PlotQA, ChartBench, MathVista, and InfographicVQA subsets. These checks establish two narrow facts: ratio-related scale errors occur under every backbone we test, and the prefer-large adapter does not transfer as a correction (InternVL2.5-8B shows only tiny differences among directions, and Qwen external subsets are unchanged across adapters). Because models, prompts, and subset sizes differ, they do not establish a common prevalence or taxonomy across models; we use them to bound the method's scope.

\section{Analysis}
\label{sec:analysis}

Scale interventions create useful alternatives; the system fails at deciding when to trust them.

\subsection{Pair Direction Is Not Symmetric}


In the development ablation (Table~\ref{tab:pair_ablation}, appendix), prefer-large pairs give the best full-validation score while prefer-small pairs reduce it by 1.51 points; yet ScaleQA-Natural shows regimes where prefer-small is best. Pair direction changes both which scale conflicts are repaired and which non-scale examples are broken. On ChartQA validation, the prefer-large adapter fixes a mean of 13/23 strict candidates across seeds while breaking 1.27\% of base-correct examples (Table~\ref{tab:seed_consistency}).

\paragraph{Fix--break asymmetry on ScaleQA-Natural.}
On ScaleQA-Natural-473 the same adapter fixes 5.71\% of examples the answer-tuned model misses but breaks 7.40\% of examples it gets right, a net loss. The group diagnostic (Table~\ref{tab:selector_group_diagnostics}) makes the reversal concrete: prefer-large policies are strong when the original prediction is too small but poor when it is already too large, and self-check is safer in the latter regime but leaves large headroom in the former. The structured verifier shows the same asymmetry at the example level despite seeing the same image and candidate list. This is why a global adapter can help the verified test-set slice yet fail as a deployment policy: the correct action depends on question form, unit evidence, and the direction of the original conflict.

\subsection{Regime Selection Is the Bottleneck}
\label{sec:selector}


The subtype reversal suggests that correction should be selected rather than applied globally. On ScaleQA-Core a simple subtype router reaches 50.71 against 49.97 for the best fixed policy, while the oracle-subtype router reaches 80.77 (Table~\ref{tab:router_results}, appendix): the missing component is regime detection, not another adapter.

\begin{table}[t]
\centering
\small
\resizebox{\linewidth}{!}{\begin{tabular}{lcc}
\toprule
Method on ScaleQA-Natural-473 (50 held-out splits) & Accuracy & Notes \\
\midrule
Self-check prompt (fixed policy) & 55.14 $\pm$ 4.30 & strongest fixed prompt \\
Label-free selector (chart-context + candidate evidence) & 56.67 $\pm$ 5.61 & paired $\Delta$ vs.\ self-check $+1.54$, 95\% CI $[+0.17,+2.86]$ \\
Label-free selector, leave-source-out & \verify{[fill]} & \verify{falls back to self-check on most held-out sources} \\
Best fixed policy chosen per split (oracle over fixed policies) & 57.14 $\pm$ 3.62 & uses split labels \\
Pairwise top-8 verifier & 54.33 & candidate comparison prompt \\
Structured-context verifier & 48.41 & image + chart context + top-4 candidates \\
Table extraction + calculator & 42.28 & DePlot-style evidence, symbolic calculation \\
\midrule
Policy-pool oracle & 90.04 & upper bound \\
Candidate top-8 oracle & 91.75 & upper bound \\
\bottomrule
\end{tabular}
}
\caption{Deployable selection versus fixed policies and oracles on ScaleQA-Natural-473. ``Label-free'' means no gold subtype or direction feature is used; it does not mean an untouched test set, so we additionally report leave-source-out evaluation. The full ladder of selector variants is in Appendix~\ref{app:selectors}.}
\label{tab:selector_main}
\end{table}

Table~\ref{tab:selector_main} compares the strongest selector with fixed policies, verifiers, and oracles. \emph{Protocol.} For each of 50 random 60/20/20 partitions, every selector configuration (feature model, stack-vs-patch mode, switch margin) and every fixed policy is scored on the 20\% dev partition only; the single best configuration is refit on train+dev and scored once on the 20\% test partition. The label-free selector reaches 56.67 against split-matched self-check at 55.14; the paired per-split difference is $+1.54$ points (bootstrap 95\% CI $[+0.19,+2.88]$; Wilcoxon $p=0.028$; 30 wins / 5 ties / 15 losses, sign test $p=0.036$). Two comparisons bound this result. A fixed policy chosen on the same dev partition scores only 52.74 on test, below self-check: naive dev-based policy choice overfits, and the selector's gain comes from switching conservatively (48.9\% of items) rather than from strong regime detection. An oracle that picks the best fixed policy per test split reaches 57.14, so the selector does not yet beat a well-chosen fixed prompt.

\begin{table}[t]
\centering
\small
\resizebox{\linewidth}{!}{\begin{tabular}{lrrrrrr}
\toprule
Held-out mining source & $N_{\text{test}}$ & $N_{\text{train}}$ & Selector (10 seeds) & Self-check & $\Delta$ & Policy oracle \\
\midrule
ChartQA train shard & 316 & 157 & 45.32 $\pm$ 4.98 & 52.22 & $-6.9$ & 83.86 \\
ChartQA val & 71 & 402 & 59.01 $\pm$ 4.22 & 64.79 & $-5.8$ & 83.10 \\
ChartQA train & 47 & 426 & 65.53 $\pm$ 7.82 & 68.09 & $-2.6$ & 80.85 \\
ChartQA test & 39 & 434 & 37.44 $\pm$ 5.82 & 38.46 & $-1.0$ & 79.49 \\
\bottomrule
\end{tabular}
}
\caption{Leave-source-out evaluation of the same selector. Each ChartQA source split that contributed ScaleQA items is held out in turn; configuration is chosen on a dev partition of the remaining sources. The selector is below self-check on every held-out source.}
\label{tab:lso}
\end{table}

The advantage does not survive source shift. Table~\ref{tab:lso} holds out each mining source in turn: the selector falls below self-check on all four, by 1.0 to 6.9 points, while the policy oracle stays at 79--84. Since selector features were also iterated over several versions on this collection (Appendix~\ref{app:selectors}), even the random-split gain is an upper estimate. The selector is therefore not a deployable method; it shows that chart-side evidence carries a weak, source-specific regime signal while most oracle headroom remains unclaimed. More aggressive selectors (statistical fusion, Bayesian priors, Bayesian-MLP) fall below self-check even on random splits (Appendix~\ref{app:selectors}).

\subsection{What the Negative Results Establish}

The negative results rule out simple explanations for the gap between generation and selection. Table extraction plus calculation tests whether the issue is missing arithmetic; it is not (42.28). Structured verifiers test whether exposing chart context and explicit scale variants is enough; it is not (48.41). Candidate-likelihood and pairwise-verifier variants test whether the model internally prefers the right candidate once shown a compact list; again no. Blind sampling (Table~\ref{tab:blind_oracle}) shows that the missing answer is usually not missing: the model generates the correct scale in \verify{79.3\%} of cases within 20 samples. Candidate generation and regime selection both matter, and the second is the larger gap.

The missing capability is calibrated regime detection: deciding when \textit{0.61} is a decimal answer, when it is a percent written as a decimal, and when a larger or smaller unit is required by the axis.

\section{Conclusion}

We studied directional scale inconsistency in chart question answering. The main empirical finding is not that one adapter universally fixes scale errors, but that correction is regime-dependent. In answer-tuned Qwen2.5-VL-7B, large-positive scale-contrastive adaptation gives a small held-out ChartQA gain and recovers nearly half of strict test-set scale errors. ScaleQA-Natural explains the boundary: self-check is stronger than a global adapter, small-positive policies help different regimes, and aggressive selectors can hurt by switching too often.

The strongest evidence is the gap between candidate availability and candidate choice. Structured verifiers, likelihood scoring, pairwise comparison, and top-8 pools often contain a correct scaled answer, yet deployable selectors fail to choose it. Our best no-leak selector improves over self-check but remains far below policy and candidate oracles. Cross-backbone and external checks show that scale-like failures transfer more readily than a correction recipe.

Overall, Chart QA systems should evaluate scale reliability by regime, with fix--break accounting and candidate-selection analysis alongside aggregate accuracy.

\clearpage
\section*{Limitations}

The strict naturally occurring scale-error subsets on official ChartQA are small. We therefore use them as recovery probes and rely on ScaleQA-Natural for broader regime analysis, but ScaleQA-Natural is still a mined diagnostic set rather than a fully independent human-authored benchmark. Its final 520 positive inclusions are author-adjudicated and independently audited by a second human reviewer, but prevalence estimates remain tied to the mining sources, and the source audit shows that the diagnostic items are not an ordinary i.i.d. benchmark split. The main positive adaptation result is obtained with Qwen2.5-VL-7B-Instruct in an answer-tuned setting. We add InternVL2.5-1B, 8B, and 26B diagnostics, including five-seed 8B adaptation, but these experiments show bounded transfer rather than universal improvement. The external PlotQA, ChartBench, MathVista, and InfographicVQA subsets are useful robustness checks, but they are not full-dataset benchmark submissions and should not be interpreted as state-of-the-art comparisons. Our stronger baselines include prompt variants, multi-candidate verification, learned and rule-based selectors, statistical policy fusion, Bayesian bucket/MLP selectors, structured chart-context verification, token-likelihood candidate scoring, and a DePlot-style table-extraction-and-calculation pipeline, but these are still lightweight baselines rather than specialized chart parsers. The scale definition focuses primarily on factor-of-100 style errors, especially percentage-decimal confusions, although ScaleQA-Natural includes smaller groups of 10x, 1000x, and unit-like candidates. Finally, the no-leak selector improves over self-check only modestly and remains far below oracle candidate selection. A deployable system would need stronger scale-regime detection using OCR, axis/unit evidence, visual features, or calibrated verifier probabilities before applying scale correction.


\bibliography{references}

\appendix
\section{Implementation Details}

\paragraph{Appendix roadmap.}
The main paper compresses the evidence into the story that scale correction is directional and that the hard part is choosing the right scale regime. This appendix gives the details behind that claim. Section~\ref{sec:appendix_data} documents how ScaleQA-Natural is mined, audited, and separated from training and official evaluation. Section~\ref{sec:appendix_tables} reports development evidence showing that the scale-adapter gain is real but narrow. Section~\ref{sec:appendix_analysis} then follows the negative results: stronger prompts, verifiers, table calculators, candidate likelihood scoring, Bayesian fusion, and cross-backbone checks often expose the right answer in the candidate pool, but still fail to choose it reliably. This is why we frame adapters and verifiers as diagnostic probes rather than as a complete repair.

\paragraph{Adapters.}
The answer-tuned setting continues from the ChartQA/TextVQA answer-only LoRA adapter used in the evaluation pipeline. The answer-only adapters use rank $r=16$, LoRA $\alpha=32$, dropout 0.05, and target the language-model projection and MLP modules \texttt{q\_proj}, \texttt{k\_proj}, \texttt{v\_proj}, \texttt{o\_proj}, \texttt{gate\_proj}, \texttt{up\_proj}, and \texttt{down\_proj}. Answer-only training uses AdamW with learning rate $2\times 10^{-5}$. The final Qwen scale adapters use large-positive pairs with $K=300$, 300 optimization steps, AdamW learning rate $10^{-6}$, weight decay 0.01, gradient clipping at 1.0, $\beta=2.0$, and margin $m=0.10$. The clean direct setting starts from the raw Qwen2.5-VL-7B-Instruct backbone and applies only the scale-contrastive update. The InternVL2.5-1B adapter check uses the same pair construction, rank $r=8$, LoRA $\alpha=16$, dropout 0.05, 200 optimization steps, and the same language-model target modules.

\paragraph{Why separate answer-tuned and clean-direct settings.}
The answer-only adapter is not a scale intervention. It mostly makes the model better aligned to short chart answers, numeric formatting, and the evaluation prompt. Keeping this setting separate matters because the clean-direct scale-only adapter is unsafe on the raw backbone (Table~\ref{tab:clean_direct}), while the answer-tuned setting can absorb a small scale-specific update without the same global collapse. This difference suggests that scale correction is not a standalone numerical operation. It depends on a model already being in a chart-answering regime where answer extraction, formatting, and short-answer calibration are reasonably stable.

\paragraph{Decoding and scoring.}
All single-adapter systems use the same short-answer prompt and deterministic short-answer decoding template: the question is followed by ``Answer briefly.'' and generation uses greedy decoding with \texttt{do\_sample=False} and \texttt{max\_new\_tokens=24}. We evaluate numeric answers with the repository evaluator, which uses normalized exact matching for text and a 5\% numerical tolerance for numeric answers.

\paragraph{ScaleQA-Natural annotation.}
Candidate scale conflicts are mined from model predictions and automatically prefilled with numeric-ratio features. The review workflow assigns labels including scale diagnostic candidate, requires image inspection, non-scale, uncertain, and duplicate. Positive labels require that the gold answer and prediction refer to the same chart quantity but differ primarily by a simple scale relation; wrong entity, wrong operation, OCR failure, arbitrary numeric miss, ambiguous question, and duplicate cases are excluded or sent to adjudication. The final ScaleQA-Natural-520 set contains author-adjudicated positive candidates after question/gold deduplication. The image-evaluable subset contains 473 examples with available ChartQA image rows. The audit workflow exports three files: the 739-row candidate audit, a conflict queue, and a 108-row final-positive adjudication queue. After this adjudication, an independent second reviewer audited the 520 final positive inclusions and agreed with all inclusion labels. The release package is intended to include item IDs, source split and row index, question, gold answer, model prediction, subtype, direction, review label, audit metadata, and scripts for reproducing the filters. We reference ChartQA images by source row rather than redistributing them.

\paragraph{PlotQA external subset.}
The PlotQA external check uses a 500-example subset sampled from the official PlotQA test parquet after filtering to examples with a single numeric answer. We use the same relaxed numeric evaluator as the ChartQA experiments. This subset is intended as an external numeric-answer transfer check rather than a full PlotQA benchmark.

\section{Data Flow and Annotation Audit}
\label{sec:appendix_data}

ScaleQA-Natural is useful only if it is not mistaken for an ordinary benchmark split or for hidden training data. Table~\ref{tab:data_flow} therefore separates five roles: answer-only tuning, scale-pair training, development analysis, diagnostic set construction, and held-out or external evaluation. The reported Qwen scale adapters are trained from ChartQA training rows and training-time prediction caches; ScaleQA-Natural is used to diagnose regimes and evaluate policy selection, not to train those scale adapters. The held-out ChartQA test result remains the official aggregate task result.

\begin{table*}[h]
\centering
\small
\resizebox{\textwidth}{!}{\begin{tabular}{p{0.18\linewidth}p{0.25\linewidth}p{0.25\linewidth}p{0.24\linewidth}}
\toprule
Component & Source & Used for & Leakage control \\
\midrule
Answer-only tuning & ChartQA/TextVQA training data & Starting adapter for the main Qwen setting & Reported separately from scale tuning; not counted as a scale-objective gain \\
Scale-pair training & ChartQA training rows and training-set prediction/log-prob caches & LoRA scale adapters and pair-mode ablations & Official ChartQA test and PlotQA are not used for scale-adapter gradients \\
Development analysis & ChartQA validation and ScaleQA-Core diagnostics & Pair-mode selection, seed checks, and failure analysis & Reported as development evidence; held-out ChartQA test is the main official result \\
ScaleQA-Natural & Mined model outputs, machine audit, and author adjudication & Diagnostic regime analysis and router experiments & Used as a diagnostic set, not as ordinary task training data for the reported Qwen scale adapters \\
Held-out/external evaluation & ChartQA test, TextVQA-500, PlotQA numeric-500 & Final task result, safety, and external robustness check & Not used to train the reported scale adapters \\
\bottomrule
\end{tabular}
}
\caption{Data flow and split hygiene.}
\label{tab:data_flow}
\end{table*}

\begin{table}[h]
\centering
\small
\resizebox{\linewidth}{!}{\begin{tabular}{llrr}
\toprule
Audit view & Source / flag & Count & Share \\
\midrule
Image-evaluable 473 & train\_shard1 & 314 & 66.4\% \\
Image-evaluable 473 & val & 71 & 15.0\% \\
Image-evaluable 473 & train & 47 & 9.9\% \\
Image-evaluable 473 & test & 39 & 8.2\% \\
Image-evaluable 473 & train\_shard2 & 1 & 0.2\% \\
Image-evaluable 473 & train\_shard3 & 1 & 0.2\% \\
Final 520 & Contains ChartQA train occurrence & 336 & 64.6\% \\
Final 520 & Contains ChartQA validation occurrence & 59 & 11.3\% \\
Final 520 & Contains ChartQA test occurrence & 23 & 4.4\% \\
\bottomrule
\end{tabular}
}
\caption{ScaleQA-Natural source audit. The set is a diagnostic slice mined from multiple model outputs and source rows, not an ordinary training/test split or leaderboard. Occurrence flags are counted over the final 520 positives; source-split rows are counted over the 473 image-evaluable items.}
\label{tab:source_audit}
\end{table}

The annotation story has two parts. First, strict factor-of-100 failures are sparse in any single ChartQA output file (Table~\ref{tab:strict_prevalence}), so a single split underestimates the phenomenon. Second, mining introduces the opposite risk: plausible numeric ratios can be false positives when the model answered the wrong entity or operation. We therefore use ScaleQA-Natural as a diagnostic set whose positives require author adjudication and an independent inclusion audit (Table~\ref{tab:annotation_audit_app}), and we keep uncertain/rejected examples visible (Table~\ref{tab:negative_audit_examples}) rather than silently discarding them. The resulting set is deliberately conservative: it favors precision of the scale-conflict label over treating every ratio-like mismatch as a scale error.

\begin{table}[h]
\centering
\small
\resizebox{\linewidth}{!}{\begin{tabular}{lrl}
\toprule
Audit item & Count/Score & Notes \\
\midrule
Candidate pool & 739 & First-pass reviewed candidates \\
Observed agreement & 0.633 & Three-way labels \\
Cohen's $\kappa$ & 0.443 & Include/review/reject \\
Conflicts for review & 271 & Exported separately \\
Final positives & 520 & Author-adjudicated \\
Second-pass include & 412 & Deterministic audit \\
Second-pass review/reject & 108 & Manually adjudicated \\
Independent inclusion audit & 520 & Second human reviewer \\
Second-review agreement & 1.000 & Final positives; zero label changes \\
\bottomrule
\end{tabular}
}
\caption{Annotation audit for ScaleQA-Natural. The deterministic metadata pass flags conflicts for author adjudication; the final positive inclusions are then independently audited by a second human reviewer.}
\label{tab:annotation_audit_app}
\end{table}

\begin{table}[h]
\centering
\small
\resizebox{\linewidth}{!}{\begin{tabular}{llrr}
\toprule
Model / pipeline & Split & Strict errors & Rate \\
\midrule
Qwen2.5-VL-7B + answer/large & ChartQA test & 10 / 2500 & 0.40\% \\
Qwen2.5-VL-3B direct/large & ChartQA test & 19 / 2500 & 0.76\% \\
Qwen2.5-VL-3B direct/large & ChartQA val & 21 / 1920 & 1.09\% \\
InternVL2.5-1B & ChartQA test & 3 / 2500 & 0.12\% \\
InternVL2.5-1B & ChartQA val & 17 / 1920 & 0.89\% \\
\bottomrule
\end{tabular}
}
\caption{Strict factor-of-100 scale-error prevalence in available ChartQA validation/test prediction files. These sparse natural subsets motivate the larger ScaleQA-Natural diagnostic set.}
\label{tab:strict_prevalence}
\end{table}

\begin{table*}[h]
\centering
\small
\resizebox{\textwidth}{!}{\begin{tabular}{lll}
\toprule
Audit label & Auto subtype & Question snippet \\
\midrule
uncertain & percent\_to\_decimal & How much Good more than Bad? \\
uncertain & value\_scaled\_up & How many diamond cutting and polishing jobs were there in the Far East i... \\
uncertain & decimal\_to\_percent & How many respondents say no? \\
uncertain & decimal\_to\_percent & How many says that internet has good influence on education? \\
uncertain & decimal\_to\_percent & How confident are people about Angela Merkel doing the right thing regar... \\
uncertain & decimal\_to\_percent & How many more people answered Very acceptable for the most and least pop... \\
\bottomrule
\end{tabular}
}
\caption{Examples from the negative/uncertain audit queue. These candidates contain a plausible scale relation but are not included as positive diagnostics without stronger evidence that the prediction and gold answer refer to the same chart quantity.}
\label{tab:negative_audit_examples}
\end{table*}

\section{Additional Results}
\label{sec:appendix_tables}

The development results explain why the paper does not present the large-positive adapter as a general correction method. On ChartQA validation, the adapter recovers strict gold-scale errors much more often than it changes ordinary examples (Tables~\ref{tab:bootstrap} and~\ref{tab:seed_consistency}), but the aggregate confidence interval is small and overlaps zero on the full split. This is exactly the diagnostic pattern we care about: the intervention finds a real failure mode, but it is not a broad accuracy engine.

\begin{table*}[h]
\centering
\small
\resizebox{\textwidth}{!}{\begin{tabular}{lrrrrrrr}
\toprule
Split & N & Base & Ours & Delta & CI low & CI high & $p(\Delta \leq 0)$ \\
\midrule
All examples & 1920 & 0.860420 & 0.862500 & 0.002080 & -0.001250 & 0.005520 & 0.113500 \\
Numeric-answer examples & 1452 & 0.860880 & 0.862810 & 0.001930 & -0.002200 & 0.006060 & 0.187800 \\
Scale-keyword numeric examples & 1179 & 0.845630 & 0.847840 & 0.002210 & -0.002880 & 0.007460 & 0.202500 \\
Strict gold-scale errors & 23 & 0.000000 & 0.113040 & 0.113040 & 0.000000 & 0.252170 & 0.036800 \\
Non-gold-scale examples & 1897 & 0.870850 & 0.871590 & 0.000740 & -0.002210 & 0.003800 & 0.327100 \\
\bottomrule
\end{tabular}
}
\caption{Bootstrap analysis on ChartQA validation for one large-positive adapter snapshot. The full-set confidence interval overlaps zero, so we emphasize the directionality and diagnostic analyses rather than overclaiming aggregate accuracy.}
\label{tab:bootstrap}
\end{table*}

\begin{table}[h]
\centering
\small
\resizebox{\linewidth}{!}{\begin{tabular}{lrrrrr}
\toprule
Split & Base & Ours & Std & Delta & N \\
\midrule
All examples & 0.860400 & 0.864000 & 0.000600 & 0.003500 & 1920 \\
Strict gold-scale errors & 0.000000 & 0.208700 & 0.036400 & 0.208700 & 23 \\
Non-gold-scale examples & 0.870800 & 0.871900 & 0.000400 & 0.001100 & 1897 \\
\bottomrule
\end{tabular}
}
\caption{ChartQA validation subset analysis for the answer-tuned large-positive scale adapter. This was used as development evidence; the main official result is the held-out test result in Table~\ref{tab:test_subset_results}.}
\label{tab:val_subset_results}
\end{table}

\begin{table*}[h]
\centering
\small
\resizebox{\textwidth}{!}{\begin{tabular}{lrrrrrrrr}
\toprule
Split & N & Base & Mean fixes & Mean breaks & Any fix & All fix & Any break & All break \\
\midrule
All examples & 1920 & 0.860400 & 0.021400 & 0.010900 & 0.005700 & 0.003100 & 0.002600 & 0.001600 \\
Base-wrong examples & 268 & 0.000000 & 0.153000 & 0.000000 & 0.041000 & 0.022400 & 0.000000 & 0.000000 \\
Base-correct examples & 1652 & 1.000000 & 0.000000 & 0.012700 & 0.000000 & 0.000000 & 0.003000 & 0.001800 \\
Strict gold-scale errors & 23 & 0.000000 & 0.565200 & 0.000000 & 0.130400 & 0.087000 & 0.000000 & 0.000000 \\
Non-gold-scale examples & 1897 & 0.870800 & 0.014800 & 0.011100 & 0.004200 & 0.002100 & 0.002600 & 0.001600 \\
\bottomrule
\end{tabular}
}
\caption{Seed-level fix and break rates on ChartQA validation.}
\label{tab:seed_consistency}
\end{table*}

\begin{table}[h]
\centering
\small
\begin{tabular}{lrrr}
\toprule
Method & Pairs per kind & Accuracy & Delta \\
\midrule
Base & 0 & 0.850700 & 0.000000 \\
Large-25 & 25 & 0.856000 & 0.005300 \\
Large-50 & 50 & 0.856000 & 0.005300 \\
Large-100 & 100 & 0.855300 & 0.004700 \\
Large-200 & 200 & 0.856000 & 0.005300 \\
Large-300 & 300 & 0.855300 & 0.004700 \\
Large-400 & 400 & 0.855300 & 0.004700 \\
Large-600 & 600 & 0.856000 & 0.005300 \\
Large-800 & 800 & 0.856000 & 0.005300 \\
\bottomrule
\end{tabular}

\caption{Large-positive pair-size curve on the ChartQA-1500 setting. Gains appear with few pairs and then saturate.}
\label{tab:size_curve}
\end{table}

Table~\ref{tab:size_curve} is another reason we avoid a ``more LoRA pairs solves scale'' claim. Gains appear quickly and then saturate, while the subtype table in the main paper shows that the opposite direction requires a different policy. The ablation in Table~\ref{tab:ablation} should therefore be read as a probe-selection experiment: pair direction changes which mistakes are repaired and which safe cases are put at risk.

\begin{table*}[t]
\centering
\small
\resizebox{\textwidth}{!}{%
\begin{tabular}{lrrrrrr}
\toprule
Method & Overall & $\Delta$ & Gold-scale & $\Delta$ & Non-scale & $\Delta$ \\
\midrule
Base & 79.00 & +0.00 & 0.00 & +0.00 & 82.81 & +0.00 \\
Scale-aware prompt & 79.00 & +0.00 & 0.00 & +0.00 & 82.81 & +0.00 \\
Direct SFT (800) & 78.20 & -0.80 & 4.35 & +4.35 & 81.76 & -1.05 \\
Direct SFT (1500) & 77.40 & -1.60 & 4.35 & +4.35 & 80.92 & -1.89 \\
Mixed contrastive & 79.40 & +0.40 & 8.70 & +8.70 & 82.81 & +0.00 \\
Balanced contrastive & 79.20 & +0.20 & 4.35 & +4.35 & 82.81 & +0.00 \\
Large-only contrastive & 79.00 & +0.00 & 4.35 & +4.35 & 82.60 & -0.21 \\
Real contrastive & 79.40 & +0.40 & 8.70 & +8.70 & 82.81 & +0.00 \\
\textbf{Real contrastive (5 seeds)} & \textbf{79.52} & \textbf{+0.52} & 11.30 & +11.30 & 82.81 & +0.00 \\
\bottomrule
\end{tabular}
}
\caption{Ablation results on ChartQA-500. Gold-scale denotes strict 100$\times$ scale-error cases where the base model predicts an answer at the wrong numerical scale.}
\label{tab:ablation}
\end{table*}


\section{Additional Transfer and Analysis Tables}
\label{sec:appendix_analysis}

This section contains the main negative evidence behind the paper's conclusion. If scale correction were a simple numeric normalization problem, stronger backbones, larger candidate pools, chart-to-table calculators, or Bayesian policy fusion should close the gap. Instead, the same pattern recurs. The correct answer is often present under some candidate or policy oracle, but deployable selectors remain directionally biased. This supports the paper's central distinction: the phenomenon transfers more robustly than the intervention.

\begin{table}[h]
\centering
\small
\resizebox{\linewidth}{!}{\begin{tabular}{llrrr}
\toprule
Backbone & Split & N & Accuracy & Strict scale errors \\
\midrule
Qwen2.5-VL-7B & ChartQA test & 2500 & 0.7804 & 10 \\
InternVL2.5-1B & ChartQA validation & 1920 & 0.7745 & 17 \\
InternVL2.5-1B & ChartQA test & 2500 & 0.7524 & 3 \\
InternVL2.5-1B & ScaleQA-Natural-520 & 473 & 0.3340 & -- \\
\bottomrule
\end{tabular}
}
\caption{Cross-backbone strict scale-like error diagnostic.}
\label{tab:cross_backbone}
\end{table}

\begin{table*}[h]
\centering
\small
\resizebox{\textwidth}{!}{\begin{tabular}{llrrrr}
\toprule
Set & Model & Base & Large & Small & Bal. \\
\midrule
ScaleQA-Natural-520 & InternVL2.5-1B & 0.3340 & 0.3446 & 0.3446 & -- \\
ScaleQA-Natural-520 & InternVL2.5-8B & 0.4672 & 0.4685 & 0.4723 & 0.4689 \\
ScaleQA-Natural-520 & InternVL2.5-26B & 0.4989 & -- & -- & -- \\
ScaleQA-stress1000 & InternVL2.5-8B & 0.7320 & 0.7272 & 0.7276 & 0.7256 \\
ScaleQA-stress1000 & Qwen2.5-VL-7B & 0.6940 & 0.6940 & 0.6940 & 0.6940 \\
PlotQA-557 & InternVL2.5-8B & 0.6912 & 0.6829 & 0.6862 & 0.6840 \\
PlotQA-500 & InternVL2.5-26B & 0.7200 & -- & -- & -- \\
ChartBench-500 & Qwen2.5-VL-7B & 0.4820 & 0.4820 & 0.4820 & 0.4820 \\
ChartBench-500 & InternVL2.5-8B & 0.3520 & 0.3512 & 0.3516 & 0.3520 \\
ChartBench-500 & InternVL2.5-26B & 0.4080 & -- & -- & -- \\
MathVista-243 & Qwen2.5-VL-7B & 0.6708 & 0.6708 & 0.6708 & 0.6708 \\
MathVista-243 & InternVL2.5-8B & 0.6337 & 0.6379 & 0.6387 & 0.6387 \\
MathVista numeric & InternVL2.5-26B & 0.6749 & -- & -- & -- \\
InfographicVQA-300 & Qwen2.5-VL-7B & 0.7767 & 0.7767 & 0.7767 & 0.7767 \\
InfographicVQA-300 & InternVL2.5-8B & 0.6067 & 0.6087 & 0.6073 & 0.6093 \\
InfographicVQA-300 & InternVL2.5-26B & 0.6300 & -- & -- & -- \\
\bottomrule
\end{tabular}
}
\caption{Full cross-backbone and external numeric transfer table.}
\label{tab:transfer_boundary_full}
\end{table*}

Tables~\ref{tab:cross_backbone}--\ref{tab:external_internvl8b} show the transfer boundary. Larger InternVL models improve the base level on ScaleQA-Natural and external numeric subsets, but the same large/small/balanced scale adapters do not become uniformly helpful. On ScaleQA-stress1000, for example, InternVL2.5-8B adapters slightly reduce accuracy relative to the base model; on MathVista and InfographicVQA the changes are small; on PlotQA the adapter can hurt. These results do not refute the existence of scale inconsistency. They show that a global scale intervention learned in one model/task setting is not the right abstraction for deployment.

\begin{table}[h]
\centering
\small
\resizebox{\linewidth}{!}{\begin{tabular}{lrrr}
\toprule
Backbone and setting & N & Accuracy & Std. \\
\midrule
InternVL2.5-1B base & 473 & 0.3340 & -- \\
InternVL2.5-1B large-positive LoRA & 473 & 0.3446 & -- \\
InternVL2.5-1B small-positive LoRA & 473 & 0.3446 & -- \\
\midrule
InternVL2.5-8B base & 473 & 0.4672 & -- \\
InternVL2.5-8B large-positive LoRA & 473 & 0.4685 & 0.0017 \\
InternVL2.5-8B small-positive LoRA & 473 & 0.4723 & 0.0035 \\
InternVL2.5-8B balanced LoRA & 473 & 0.4689 & 0.0025 \\
\bottomrule
\end{tabular}
}
\caption{InternVL scale adaptation on ScaleQA-Natural-520.}
\label{tab:internvl_adapter}
\end{table}

\begin{table}[h]
\centering
\small
\resizebox{\linewidth}{!}{\begin{tabular}{lrrr}
\toprule
InternVL2.5-8B setting & N & Accuracy & Std. \\
\midrule
Base & 1000 & 0.7320 & -- \\
Large-positive adapter & 1000 & 0.7272 & 0.0016 \\
Small-positive adapter & 1000 & 0.7276 & 0.0034 \\
Balanced adapter & 1000 & 0.7256 & 0.0023 \\
\bottomrule
\end{tabular}
}
\caption{InternVL2.5-8B on a synthetic ScaleQA-stress1000 set.}
\label{tab:internvl8b_stress}
\end{table}

\begin{table}[h]
\centering
\small
\begin{tabular}{lr}
\toprule
Prompting strategy & Accuracy \\
\midrule
Direct answer & 0.5243 \\
Axis-aware prompt & 0.5307 \\
Math-check prompt & 0.5137 \\
Scale prompt & 0.4905 \\
Self-check prompt & 0.5455 \\
Unit-guard prompt & 0.5180 \\
Majority vote & 0.5264 \\
Oracle prompt choice & 0.6744 \\
\bottomrule
\end{tabular}

\caption{Training-free prompting baselines on ScaleQA-Natural-520.}
\label{tab:tta_baselines_app}
\end{table}

\begin{table}[h]
\centering
\small
\resizebox{\linewidth}{!}{\begin{tabular}{lccc}
\toprule
Method & Overall & Pred. too large & Pred. too small \\
\midrule
Candidate majority & 0.5159 & 0.3029 & 0.7371 \\
Free-form verifier & 0.4588 & 0.1411 & 0.7888 \\
Letter verifier & 0.4651 & 0.1411 & 0.8017 \\
Scale-guard letter verifier & 0.4609 & 0.1452 & 0.7888 \\
Candidate oracle & 0.8288 & 0.7593 & 0.9009 \\
\bottomrule
\end{tabular}
}
\caption{Multi-candidate verifier on ScaleQA-Natural-520. Candidate oracle selects the best answer from the same candidate pool and is an upper bound.}
\label{tab:candidate_verifier}
\end{table}

The selector and verifier results give the sharpest version of the bottleneck. Self-check prompting is already a strong fixed policy, and the best no-leak selector improves on it only modestly (Table~\ref{tab:deployable_selectors}). Yet the candidate and policy oracles are much higher. The gap is not caused by a failure to generate alternatives: it is caused by failures to rank alternatives under chart-specific scale evidence. The structured verifier makes this especially clear. It is strong when the prediction is too small but nearly collapses when the prediction is too large (Table~\ref{tab:structured_verifier_direction}), even though the augmented oracle is high in both directions.

\begin{table*}[h]
\centering
\small
\resizebox{\textwidth}{!}{\begin{tabular}{lcc}
\toprule
Baseline or selector & Accuracy & Notes \\
\midrule
Self-check prompt & 0.5514 & Best fixed prompt on held-out splits \\
v8 no-leak selector & 0.5432 $\pm$ 0.0501 & Question, answer, numeric, and chart-context features \\
v9 expanded selector & 0.5501 $\pm$ 0.0460 & Adds structured/pairwise policies but does not improve \\
v10b no-leak selector & 0.5667 $\pm$ 0.0561 & Best clean selector; adds corrected TTA-v2 evidence \\
v11 log-prob selector & 0.5667 $\pm$ 0.0561 & Log-prob evidence preserves but does not improve v10b \\
v13 rationale selector & 0.5667 $\pm$ 0.0561 & Rationale-verifier evidence preserves v10b \\
v14/v15 candidate-likelihood selectors & 0.5667 $\pm$ 0.0561 & Answer/large/base likelihood evidence preserves v10b \\
v16/v17 top-8 selectors & 0.5667 $\pm$ 0.0561 & Larger candidate pools preserve v10b but raise only oracle \\
Statistical policy fusion & 0.5309 $\pm$ 0.0470 & Beta/Wilson reliability fusion \\
Bucketed Bayesian selector & 0.5239 $\pm$ 0.0480 & Bucketed policy priors; over-switches \\
Bayesian-MLP selector & 0.5218 $\pm$ 0.0540 & Bucket priors plus hidden-layer selector; over-switches \\
Pairwise top-8 verifier & 0.5433 & Candidate comparison prompt \\
Candidate likelihood top-8 & 0.4503 & Log-probability candidate scoring \\
Structured-context verifier & 0.4841 & Image, chart context, and top-4 candidates \\
Structured scale-audit verifier & 0.4376 & Explicit unit/scale audit prompt \\
Table extraction + calculator & 0.4228 & DePlot-style chart-to-table evidence, then symbolic calculation \\
\midrule
Original policy-pool oracle & 0.7023 $\pm$ 0.0384 & Upper bound over early policy pool \\
Expanded policy-pool oracle & 0.9004 & Upper bound after adding prompts/verifiers \\
Candidate top-8 oracle & 0.9175 & Correct answer often appears in candidate set \\
Augmented structured oracle & 0.9260 & Upper bound after adding scaled candidate variants \\
\bottomrule
\end{tabular}
}
\caption{Deployable selectors and strong baselines on ScaleQA-Natural. Learned and rule-based selectors use held-out splits, while fixed prompting, verifier, and table-calculator rows are full-set scores. The gap to the oracle rows shows that policy selection remains the bottleneck.}
\label{tab:deployable_selectors}
\end{table*}

\begin{table}[h]
\centering
\small
\resizebox{\linewidth}{!}{\begin{tabular}{lccc}
\toprule
Quantity & Mean & Std. & Additional statistic \\
\midrule
No-leak selector & 0.5667 & 0.0561 & 50 held-out splits \\
Self-check prompt & 0.5514 & 0.0430 & same splits \\
Selector - self-check & 0.0154 & 0.0492 & 95\% CI [0.0017, 0.0286] \\
Best fixed policy & 0.5714 & 0.0362 & split oracle over fixed policies \\
Policy oracle & 0.9004 & 0.0238 & upper bound \\
Switch rate & 48.9\% & 39.3\% & not using self-check \\
\bottomrule
\end{tabular}
}
\caption{No-leak selector stability on ScaleQA-Natural held-out splits. The improvement over self-check is modest but positive under a split-level bootstrap; the oracle gap remains much larger.}
\label{tab:selector_stability}
\end{table}

\begin{table*}[h]
\centering
\small
\resizebox{\textwidth}{!}{\begin{tabular}{lrrrrrr}
\toprule
Group & N & Selector & Ans. & Large & Self-check & Best fixed \\
\midrule
Direction: prediction too large & 241 & 0.4860 & 0.2946 & 0.1452 & 0.4357 & 0.4357 \\
Direction: prediction too small & 232 & 0.6534 & 0.7328 & 0.8621 & 0.6595 & 0.8621 \\
Source: test & 39 & 0.4275 & 0.3590 & 0.5385 & 0.3846 & 0.5385 \\
Source: val & 71 & 0.6250 & 0.5915 & 0.6338 & 0.6479 & 0.6479 \\
Source: train & 47 & 0.6303 & 0.6170 & 0.5319 & 0.6809 & 0.6809 \\
Source: train shard1 & 314 & 0.5587 & 0.4904 & 0.4522 & 0.5191 & 0.5191 \\
Subtype: decimal to percent & 109 & 0.5193 & 0.2752 & 0.1468 & 0.3853 & 0.3853 \\
Subtype: percent to decimal & 74 & 0.6411 & 0.8243 & 0.9324 & 0.7297 & 0.9324 \\
Subtype: value scaled down & 118 & 0.6578 & 0.7034 & 0.8729 & 0.6102 & 0.8729 \\
Subtype: value scaled up & 102 & 0.4098 & 0.2549 & 0.0784 & 0.4216 & 0.4216 \\
\bottomrule
\end{tabular}
}
\caption{Selector and fixed-policy behavior by source, direction, and subtype. ``Selector'' is the v10b held-out selector score averaged over test-fold appearances. Fixed-policy columns are item-level diagnostics from the candidate pool and are included to show directional reversal, not as a new benchmark leaderboard.}
\label{tab:selector_group_diagnostics}
\end{table*}

\begin{table}[h]
\centering
\small
\resizebox{\linewidth}{!}{\begin{tabular}{lccc}
\toprule
Method & Overall & Pred. too large & Pred. too small \\
\midrule
Structured-context verifier & 0.4503 & 0.0954 & 0.8190 \\
Structured scale-audit verifier & 0.4376 & 0.0996 & 0.7888 \\
Augmented structured oracle & 0.9260 & 0.9419 & 0.9095 \\
\bottomrule
\end{tabular}
}
\caption{Structured verifier by scale direction. Adding chart context and scaled candidate variants makes the oracle high, but deployable verifiers remain directionally biased.}
\label{tab:structured_verifier_direction}
\end{table}

\begin{table}[h]
\centering
\small
\resizebox{\linewidth}{!}{\begin{tabular}{lrrr}
\toprule
Subset & Reason & Calculator & Either \\
\midrule
All examples & 0.3890 & 0.4080 & 0.4228 \\
Percent-to-decimal & 0.8378 & 0.8243 & 0.8378 \\
Value scaled down & 0.6949 & 0.8220 & 0.8475 \\
Decimal-to-percent & 0.0275 & 0.0183 & 0.0183 \\
Value scaled up & 0.0098 & 0.0098 & 0.0098 \\
\bottomrule
\end{tabular}
}
\caption{DePlot-style table-extraction plus calculator baseline by subtype on ScaleQA-Natural. The method is strong for some regimes but fails on the opposite direction, again showing that scale correction is regime-dependent.}
\label{tab:table_calculator_subtypes}
\end{table}

The table-calculator baseline is a useful stress test because it is closer to symbolic reasoning than to answer-style prompting. It performs well on percent-to-decimal and value-scaled-down cases, but almost never solves decimal-to-percent or value-scaled-up cases (Table~\ref{tab:table_calculator_subtypes}). That asymmetry mirrors the adapter story: changing the evidence source does not remove the need to identify the regime first.

\begin{table}[h]
\centering
\small
\resizebox{\linewidth}{!}{\begin{tabular}{lccccc}
\toprule
Policy pool & Best fixed & Learned & +Chart ctx. & Subtype oracle & Inst. oracle \\
\midrule
Prompt regimes (473) & 0.5408 & 0.5366 & 0.5358 & 0.6192 & 0.6701 \\
Adapters+prompts (403) & 0.5032 & 0.4467 & 0.4615 & 0.8010 & 0.8410 \\
\bottomrule
\end{tabular}
}
\caption{Learned policy routing with question/base-answer numeric features and optional chart-side context extracted from the image. Values are 50-seed means.}
\label{tab:chart_context_router}
\end{table}

\begin{table}[h]
\centering
\small
\resizebox{\linewidth}{!}{\begin{tabular}{llr}
\toprule
Backbone & Method & Accuracy \\
\midrule
Qwen2.5-VL-7B & Answer-tuned & 0.902 \\
Qwen2.5-VL-7B & Answer-tuned + large-positive scale & 0.898 \\
Qwen2.5-VL-7B & Base + scale prompt & 0.832 \\
Qwen2.5-VL-7B & TTA majority & 0.892 \\
Qwen2.5-VL-7B & TTA oracle prompt choice & 0.904 \\
InternVL2.5-1B & Base & 0.600 \\
\bottomrule
\end{tabular}
}
\caption{External numeric PlotQA-500 check.}
\label{tab:plotqa_external}
\end{table}

\begin{table}[h]
\centering
\small
\resizebox{\linewidth}{!}{\begin{tabular}{lrrrrr}
\toprule
Dataset & N & Base & Large & Small & Balanced \\
\midrule
PlotQA & 557 & 0.6912 & 0.6829 & 0.6862 & 0.6840 \\
ChartBench & 500 & 0.3520 & 0.3512 & 0.3516 & 0.3520 \\
MathVista & 243 & 0.6337 & 0.6379 & 0.6387 & 0.6387 \\
InfographicVQA & 300 & 0.6067 & 0.6087 & 0.6073 & 0.6093 \\
\bottomrule
\end{tabular}
}
\caption{External numeric checks with InternVL2.5-8B.}
\label{tab:external_internvl8b}
\end{table}

\begin{table}[h]
\centering
\small
\resizebox{\linewidth}{!}{\begin{tabular}{lrrrrl}
\toprule
Subtype & N & Base & Large-pos. & Small-pos. & Best direction \\
\midrule
Decimal to percent & 109 & 0.0917 & 0.1560 & 0.7890 & Small \\
Value scaled up & 102 & 0.0784 & 0.0784 & 0.7059 & Small \\
Percent to decimal & 74 & 0.9054 & 0.9324 & 0.0270 & Large \\
Value scaled down & 118 & 0.8475 & 0.8559 & 0.0593 & Large \\
\bottomrule
\end{tabular}
}
\caption{ScaleQA-Natural subtype analysis on the 403-example reviewed set.}
\label{tab:subtype_results}
\end{table}

The subtype and taxonomy tables explain why the task is hard for deployable selectors. Most image-evaluable examples contain explicit percent evidence, but the correct action is not always ``add a percent'' or ``multiply by 100''. In decimal-to-percent and value-scaled-up cases, the small-positive policy is often best; in percent-to-decimal and value-scaled-down cases, the large-positive policy is best (Table~\ref{tab:subtype_results}). The same words---percent, average, difference, value lookup---appear across regimes (Table~\ref{tab:fine_taxonomy}), so a selector needs to reason over the chart, the question, the base answer, and the candidate answer jointly rather than keying off a single unit token.

\begin{table}[h]
\centering
\small
\resizebox{\linewidth}{!}{\begin{tabular}{lrr}
\toprule
Group & Count & Share \\
\midrule
\multicolumn{3}{l}{\textit{Question pattern}} \\
Percent / percentage & 133 & 28.1 \\
Average / mean & 75 & 15.9 \\
Value lookup & 75 & 15.9 \\
Difference / comparison & 65 & 13.7 \\
Maximum / highest & 44 & 9.3 \\
Minimum / lowest & 30 & 6.3 \\
\midrule
\multicolumn{3}{l}{\textit{Unit evidence}} \\
Percent in axis or question & 439 & 92.8 \\
No explicit unit & 14 & 3.0 \\
Large-unit word & 13 & 2.7 \\
Ratio or rate & 6 & 1.3 \\
Currency & 1 & 0.2 \\
\bottomrule
\end{tabular}
}
\caption{Fine-grained taxonomy of the image-evaluable ScaleQA-Natural subset. Counts summarize question patterns and unit evidence; categories are not mutually exclusive across the two blocks.}
\label{tab:fine_taxonomy}
\end{table}

\begin{table}[h]
\centering
\small
\resizebox{\linewidth}{!}{\begin{tabular}{lrrrrrr}
\toprule
Evaluation set & N & Base & Direct-large & Delta & Direct-small & Delta \\
\midrule
ChartQA validation & 1920 & 0.7536 & 0.7219 & -0.0318 & 0.6948 & -0.0589 \\
ChartQA test & 2500 & 0.7804 & 0.7456 & -0.0348 & 0.7136 & -0.0668 \\
ScaleQA-Core & 403 & 0.4591 & 0.4367 & -0.0223 & 0.4243 & -0.0347 \\
InternVL-mined & 192 & 0.1510 & 0.0938 & -0.0573 & 0.3542 & +0.2031 \\
\bottomrule
\end{tabular}
}
\caption{Clean direct scale-only adaptation from the raw Qwen2.5-VL-7B backbone.}
\label{tab:clean_direct}
\end{table}

Finally, Table~\ref{tab:clean_direct} is the safety counterexample. A scale-only update from the raw Qwen backbone degrades ChartQA validation, ChartQA test, and ScaleQA-Core, even though it helps a separate InternVL-mined subset in the opposite direction. This is why the paper treats scale adapters as probes of directional asymmetry, not as a universal repair. The case studies in Table~\ref{tab:case_study} give the qualitative version of the same lesson: the model can repair 0.61$\rightarrow$61 and still over-correct a true small decimal such as 0.0593.

\begin{table*}[h]
\centering
\small
\begin{tabular}{lp{0.47\linewidth}lll}
\toprule
Type & Question & Gold & Base & Calibrated \\
\midrule
Recovered & How much times Madagascar greater than Central Europe data? & 32 & 0.032 & 32 \\
Recovered & What the difference in value between Asia and Caribbean? & 1.18 & 0.18 & 1.18 \\
Recovered & What is the highest value of gray bar? & 61 & 0.61 & 61 \\
Recovered & What is the difference in \% between Poland and Australia for people who are worried about losing their job? & 17.05 & 0.1705 & 17.05 \\
Over-correction & Find out the average of the bottom three countries ?? & 0.0593 & 0.06 & 6.03333 \\
\bottomrule
\end{tabular}

\caption{Representative recoveries and an over-correction failure.}
\label{tab:case_study}
\end{table*}

\paragraph{Overall appendix takeaway.}
Taken together, the appendix strengthens the main claim in a deliberately conservative way. Scale-like mistakes are not frequent enough in a single benchmark split to support sweeping accuracy claims, but they are systematic enough to reveal a reliability failure once mined and audited. Global adapters expose the directionality, training-free self-checks provide a strong practical baseline, and expanded candidate pools show that many failures are recoverable in principle. The unresolved problem is selecting the right regime without gold subtype labels. This is the reason the paper argues for regime-aware scale reliability rather than presenting unconditional scale correction as solved.


\end{document}
